\documentclass[10pt,a4paper]{amsart}
\usepackage[margin=19mm]{geometry}
\usepackage{amsmath,amssymb,mathtools}
\usepackage[scr=rsfs,cal=euler]{mathalfa}
\usepackage{xfrac}
\usepackage[unicode,hidelinks,bookmarks=false]{hyperref}
\newcommand{\ie}{i.e.\ }
\newcommand{\cf}{cf.\ }
\newcommand{\resp}{respectively\ }
\newcommand{\Subsection}{\subsection}
\providecommand{\nobreakdash}{\nobreak\hbox{-}}
\setlength{\emergencystretch}{2em}
\multlinegap0pt
\usepackage{amssymb}
\usepackage[shortlabels]{enumitem}
\usepackage{mathtools}
\usepackage{algorithm}\usepackage{algpseudocode}
\newtheorem{theorem}{Theorem}
\newtheorem{lemma}{Lemma}
\newcommand{\HSOS}{\widehat{\mathbf{\Sigma}}}
\newcommand{\SOS}{\mathbf{\Sigma}}
\newcommand{\defeq}{\coloneq} 
\newcommand{\mbS}{\mathbb{S}}
\newcommand{\reals}{\mathbb{R}}
\title{Disjunctive Sum of Squares}
\author{Amir Ali Ahmadi, Sanjeeb Dash, Yixuan Hua, Bartolomeo Stellato}
\date{Source: arXiv:2605.28674v1 (CC BY 4.0)}
\begin{document}
\maketitle

\noindent\emph{Source and licence.} Disjunctive Sum of Squares. Version arXiv:2605.28674v1 (CC BY 4.0), distributed by arXiv under the Creative Commons Attribution 4.0 International licence, http://creativecommons.org/licenses/by/4.0/, as recorded in the arXiv licence metadata for this exact version and shown on the abstract page https://arxiv.org/abs/2605.28674v1. Full licence notice: \texttt{licenses/arxiv-2605.28674-CC-BY-4.0.txt}.\par\smallskip

\noindent\textit{Editorial scope.} Counted unit: the article's theorem thm:degd\_homo (source0.tex lines 852--869). The article's complete original proof of it is delivered with it (source0.tex lines 987--1043, one \texttt{proof} environment of 57 lines, which the article places away from the statement). No mathematics is written, repaired or re-derived: the statement and the proof are the article's own lines, in the article's own notation, and no retained line is substituted or re-flowed.  Retained uncounted as the source-local context the proof needs: lines 909--921; lines 913--915.  Nothing was omitted from any retained span, so the package carries no omission record.  No retained line contains a citation, so the wrapper carries no bibliography and no bibliography entry.  No retained line contains a figure environment, an image-inclusion command or drawing code, so no image is delivered and the package declares no asset dependency.  Editorial formatting only, in addition: an AMS \texttt{amsart} preamble that restates the article's own declarations and macros used by the retained lines, namely the environments \texttt{theorem}, \texttt{lemma} on separate counters, together with the standard packages those lines need.  The retained environments print in this document as Theorem 1, Lemma 1 in order of appearance, whereas the article prints the same items with the numbering of its own full document.  The complete native source of the article as distributed in the arXiv e-print for this exact version is bundled unchanged as \texttt{sources/201555/source0.tex}.
\begin{lemma}\label{thm:degd_nonhomo}
Consider a (not necessarily homogeneous) polynomial~$p$ in $n$ variables and degree at most $d$, where $d$ is even. If $p(x^*)>0$ for some $x^*\in\reals^n$, then for $t>0$ sufficiently large, $p(x)$ can be written as
\[p(x)=\left(\frac{p(x^*)}{2}-tg_d(x-x^*)\right)+\sigma(x),\]
where $\sigma(x)\in\SOS_{n,d}$ and 
\begin{equation}\label{eq:degd_nonhomo_g}
  g_d(x)\defeq\sum_{k=1}^{d/2}\|x\|_2^{2k}.  
\end{equation}
In particular, the above statement holds for any 
% $t$ such that
\[t\ge \tilde{T}(p;x^*)\defeq\dfrac{n}{2p(x^*)}\|q\|^2_{\infty} + \|q\|_1,\]
where $q(x)=p(x+x^*)-p(x^*)$. 
% For any multi-index $\alpha$, $q_{\alpha}$ denotes the coefficient of the term $x^{\alpha}$ in $q(x)$, and $\|q\|_1$ denotes the sum of the absolute value of the coefficients of $q$.
\end{lemma}

\begin{equation}
  g_d(x)\defeq\sum_{k=1}^{d/2}\|x\|_2^{2k}.  
\end{equation}

\begin{theorem}\label{thm:degd_homo}
For a polynomial $p\in H_{n,d}$ with $d$ even, if $p(x^*)>0$ for some $x^*\in \mbS^{n-1}$, then for $t>0$ sufficiently large, $p(x)$ can be written as
\begin{equation}\label{eq:degd_homo}
    p(x)=\left(\frac{p(x^*)}{2}(x^Tx^*)^d-th_d(x;x^*)\right)+\sigma(x),
\end{equation}
where $\sigma(x)\in\HSOS_{n,d}$ and 
\begin{equation}\label{eq:degd_homo_h}
h_d(x;x^*)\defeq\sum_{k=1}^{d/2}\left(\|x\|_2^2-(x^Tx^*)^2\right)^k(x^Tx^*)^{d-2k}.
\end{equation}
%TEMP:
%$$c(x)=\frac{p(x^*)}{2}(x^Tx^*)^d-t\left( \sum_{k=1}^{d/2}\left(\|x\|_2^2-(x^Tx^*)^2\right)^k(x^Tx^*)^{d-2k} \right)$$
%
In particular, the above statement holds for any
\begin{equation}\label{eq:degd_homo_t}
t\ge T(p;x^*)\defeq\sup_{U\in\reals^{n\times n}:U^TU=I_n}\left( \frac{n}{2p(x^*)}\|p(Ux)\|_{\infty}^2 + \|p(Ux)\|_1\right).  
\end{equation}
% Here, $\|p\|_{\infty}$ denotes the largest coefficient of~$p$ in absolute value, and $\|p\|_1$ denotes the sum of the absolute values of the coefficients of~$p$. 
\end{theorem}

\begin{proof}[Proof of Theorem~\ref{thm:degd_homo}]
We first consider the special case where $x^*=e_1$. Let $z\defeq(x_2,\ldots,x_n)$. We dehomogenize~$p$ by defining $\tilde{p}(z)\defeq p(1,z)$. Then $\tilde{p}(z)$ is a polynomial with $\tilde{p}(0)=p(e_1)>0$. 
By Lemma~\ref{thm:degd_nonhomo}, whenever $t\ge \tilde{T}(\tilde{p};0)$, $\tilde{p}(z)$ admits the representation
\[\tilde{p}(z)=\left(\frac{p(e_1)}{2}-tg_d(z)\right)+\tilde{\sigma}(z),\]
where $\tilde{\sigma}(z)\in\SOS_{n-1,d}$ and $g_d$ is defined as in~\eqref{eq:degd_nonhomo_g}. 
% \[g_d(z)=\sum\limits_{k=1}^{d/2}\|z\|_2^{2k}.\]
Next, homogenizing $\tilde{p}(z)$ yields the identity:
\[\begin{aligned}
p(x)&=p(x_1,z)=x_1^d \tilde{p}(z/x_1)\\
&=\left(\frac{p(e_1)}{2}x_1^d-t\sum\limits_{k=1}^{d/2}\|z\|_2^{2k}x_1^{d-2k}\right)+x_1^d \tilde{\sigma}(z/x_1)\\
&=\left(\frac{p(e_1)}{2}x_1^d-t\sum\limits_{k=1}^{d/2}(\|x\|_2^2-x_1^2)^{k}x_1^{d-2k}\right)+\sigma(x)\\
&=\left(\frac{p(e_1)}{2}x_1^d-th_d(x;e_1)\right)+\sigma(x),
\end{aligned}\]
where $\sigma(x)=x_1^d \tilde{\sigma}(z/x_1)\in\HSOS_{n,d}$, since homogenization preserves the sos property. Note that
\[\begin{aligned}
T(p;e_1)&\ge \dfrac{n}{2p(e_1)}\|p\|_{\infty}^2+\|p\|_1\\
&= \dfrac{n}{2\tilde{p}(0)}\|\tilde{p}\|^2_{\infty} + \|\tilde{p}\|_1\\
&=\tilde{T}(\tilde{p};0),
\end{aligned}\]
% Hence, whenever $t\ge T(p;e_1)$, we have 
% \[\begin{aligned}
% t&\ge \dfrac{n}{2p(e_1)}\|p\|_{\infty}^2+\|p\|_1\\
% &= \dfrac{n}{2\tilde{p}(0)}\|\tilde{p}\|^2_{\infty} + \|\tilde{p}\|_1\\
% &=\tilde{T}(\tilde{p};0),
% \end{aligned}\]
where the second line follows from the fact that the coefficients of~$p$ remain the same after dehomogenization. 
Therefore, in view of the fact that $h_d(x;e_1)$ is sos, the claim of Theorem~\ref{thm:degd_homo} follows in the case where $x^*=e_1$.

% we derive $\hat{T}_d(p;e_1)$ by upper bounding $\tilde{T}_d(\hat{p};e_1)$. 
% Note that the coefficients of $p(x)$ and $\hat{p}(z)$ are identical, we have
% \[
% T(\hat{p})=\dfrac{n}{2A}\max\limits_{1\le i\le n}\hat{p}_i^2 + \sum\limits_{\alpha\neq 0}|\hat{p}_{\alpha}|
% \le \dfrac{n}{2A}\|p\|_{\infty}^2+\|p\|_1= \tilde{T}(p;e_1),
% \]
% \[\begin{aligned}
% \tilde{T}_d(\hat{p};e_1)&=\dfrac{n}{2A}\max\limits_{1\le i\le n}\hat{p}_i^2 + \|\hat{p}\|_1\\
% &\le \dfrac{n}{2A}\|p\|_{\infty}^2+\|p\|_1= \hat{T}_d(p;e_1),
% \end{aligned}\]
% which completes the proof.
% \end{proof}

% \begin{proof}
To establish the claim for a general point $x^*\in \mbS^{n-1}$, observe that for any such point there exists a matrix $U\in\reals^{n\times n}$ with $U^TU=I$, such that~$Ue_1=x^*$. Define $p_U(x)\defeq p(Ux)$. Then $p_U$ is a polynomial with $p_U(e_1)=p(x^*)>0$. Applying the statement of the theorem to~$p_U$, whenever $t\ge T(p_U;e_1)$, we have
\[
p_U(x)=\left(\frac{p_U(e_1)}{2}(x^Te_1)^d-t\sum\limits_{k=1}^{d/2}(\|x\|_2^2-(x^Te_1)^2)^{k}(x^Te_1)^{d-2k}\right)+\hat{\sigma}(x),\]
where $\hat{\sigma}(x)\in\HSOS_{n,d}$.
Substituting $x$ with $U^Tx$ gives
\[\begin{aligned}
p(x)&=p_U(U^Tx)\\
&=\left(\frac{p(x^*)}{2}(x^TUe_1)^d-t\sum\limits_{k=1}^{d/2}(\|U^Tx\|_2^2-(x^TUe_1)^2)^{k}(x^TU  e_1)^{d-2k}\right)+\hat{\sigma}(U^Tx)\\
&=\left(\frac{p(x^*)}{2}(x^Tx^*)^d-t\sum\limits_{k=1}^{d/2}(\|x\|_2^2-(x^Tx^*)^2)^{k}(x^Tx^*)^{d-2k}\right)+\hat{\sigma}(U^Tx)\\
&=\left(\frac{p(x^*)}{2}(x^Tx^*)^d-th_d(x;x^*)\right)+\hat{\sigma}(U^Tx),
\end{aligned}
\]
where in the second line we use the fact that the 2-norm is unitarily invariant. 
Since ${\hat{\sigma}(x)\in\HSOS_{n,d}}$ and the sos property is preserved under a linear change of variables, we have ${\hat{\sigma}(U^Tx)\in\HSOS_{n,d}}$. Finally, note that $T(p_U;e_1)=T(p;x^*)$, since the set of orthogonal matrices is invariant under multiplication by any of its elements.
\end{proof}

\end{document}
